% !TeX root = ./main.tex

% 双盲评审开关：送盲审时保持 \sysublindtrue（隐去姓名、学号、导师及致谢敏感信息）；
% 答辩后提交终稿归档时改为 \sysublindfalse 并在对应分支填入实名信息。
\newif\ifsysublind
\sysublindtrue

\ifsysublind
\sysusetup{
  title                = {面向时限服务的无人机集群\newline 协同轨迹规划与计算任务分配方法研究},
  title*               = {Cooperative Trajectory Planning and Computation Task Allocation for UAV Swarms Providing Deadline-Constrained Services},
  author               = {},
  author*              = {},
  % 注：中山大学航空航天学院专业学位硕士招生目录专业名称为“机械”（085500）；送审前请与教务系统核对是否需标注方向
  speciality           = {机械},
  speciality*          = {Mechanical},
  professional-type    = {专业学位},
  professional-type*   = {Professional Degree},
  supervisor           = {},
  supervisor*          = {},
  department           = {航空航天学院},
  student-id           = {},
  math-font            = xits,
  cite-style           = super,
}
\else
\sysusetup{
  title                = {面向时限服务的无人机集群\newline 协同轨迹规划与计算任务分配方法研究},
  title*               = {Cooperative Trajectory Planning and Computation Task Allocation for UAV Swarms Providing Deadline-Constrained Services},
  author               = {学位申请人姓名},
  author*              = {Author Name},
  speciality           = {机械},
  speciality*          = {Mechanical},
  professional-type    = {专业学位},
  professional-type*   = {Professional Degree},
  supervisor           = {指导教师姓名~教授},
  supervisor*          = {Prof. Supervisor Name},
  department           = {航空航天学院},
  student-id           = {23000000},
  math-font            = xits,
  cite-style           = super,
}
\fi

% 扉页“学位”一行显示的内容（对应 degree-type=professional 专业学位硕士）
\newcommand\thesisdegreename{专业硕士}

% 自动检测我们定制的 Dev Container 环境并替换 CJK 字体，保持跨平台兼容性
\IfFileExists{/.sysu-thesis.devcontainer}{
  \sysusetup{cjk-font = noto}
}{}

% 图表公式编号改为“章-序号”（图2-1、表2-1、式(2-1)），对齐航空航天学院 Word 模板
\renewcommand{\thefigure}{\thechapter-\arabic{figure}}
\renewcommand{\thetable}{\arabic{chapter}-\arabic{table}}
\renewcommand{\theequation}{\arabic{chapter}-\arabic{equation}}

% 加载宏包

% 定理类环境宏包
\usepackage{amsthm}

% 中英文双标题
\usepackage{bicaption}

% 表注
\usepackage{threeparttable}

% 跨页表格
\usepackage{longtable}

% 跨页自动换行表格（cls 已加载 tabularx，此处仅定义左对齐列类型）
\newcolumntype{Y}{>{\raggedright\arraybackslash\hspace{0pt}}X}

% 算法
\usepackage[ruled,linesnumbered]{algorithm2e}
%\renewcommand{\thealgocf}{\arabic{chapter}-\arabic{algocf}}% 将算法标题由x.x改为x-x，慎用

% 代码块
\usepackage{listings}

% SI 量和单位
\usepackage{siunitx}

% 参考文献使用 BibTeX + natbib 宏包
% 顺序编码制
\usepackage[sort]{natbib}
\bibliographystyle{sysuthesis-numerical}

% 参考文献使用 BibLaTeX 宏包
% \usepackage[style=sysuthesis-numeric]{biblatex}
% 声明 BibLaTeX 的数据库
% \addbibresource{reference.bib}

% 配置图片的默认目录
% \graphicspath{{image/}}

% 数学命令
\makeatletter
\newcommand\dif{%  % 微分符号
  \mathop{}\!%
  \ifsysu@math@style@TeX
    d%
  \else
    \mathrm{d}%
  \fi
}
\makeatother
\newcommand\eu{{\symup{e}}}
\newcommand\iu{{\symup{i}}}

% 用于写文档的命令
\DeclareRobustCommand\cs[1]{\texttt{\char`\\#1}}
\DeclareRobustCommand\env[1]{\texttt{#1}}
\DeclareRobustCommand\pkg[1]{\textsf{#1}}
\DeclareRobustCommand\file[1]{\nolinkurl{#1}}

% hyperref 宏包在最后调用
\usepackage{hyperref}
% 封面语义断行只作用于排版，PDF 元数据保留连续的完整题名。
\pdfstringdefDisableCommands{\def\newline{}}
